\renewcommand{\stagefigureprefix}{1}\setcounter{figure}{0}

\reporttitle{第一阶段：环境验证与向量执行过程}

本阶段记录香山仿真环境部署、hello/vadd运行与向量执行调试。2026年10月9日记录展示环境建立、启动差分修复和原波形过程；10月25日的部署复核使用既有硬件构建产物，运行配置个人标识的英文hello程序，并完成含时钟信号的FST调试链。两组实验分别按各自源码、输入镜像与事件坐标说明。\par

\section{验证环境与协同仿真框架}

本阶段建立香山RTL仿真器、NEMU参考模型与裸机程序协同工作的运行入口。hello程序用于检查启动流程和串口输出，vadd程序用于观察向量配置、数据加载、按元素运算、写回与提交。通过共同启动镜像和差分比较，将工具链生成的产物与实际运行结果连接起来，为后续指令扩展提供可复用的环境。\par

\subsection{构建链路与运行产物}

协同仿真由硬件执行端、软件参考端和目标程序共同组成。香山EMU负责模拟硬件的执行过程，NEMU reference负责按指令语义推进参考状态，裸机程序提供可重复的输入并检查软件可见结果。共同flash连接复位后的启动过程与程序入口。这几类产物在一次运行中同时参与工作，部署过程因此围绕它们之间的接口组织。\par


硬件侧首先通过Java与Mill解析Scala、Chisel工程及其依赖，再由firtool处理硬件中间表示，供Verilator生成宿主侧仿真代码。宿主C++编译器将这些代码构建为EMU。每一步的产物都会成为下一步的输入：Scala工程配置影响生成的硬件，硬件生成结果影响仿真代码，宿主工具链则决定最终可执行文件的构建方式。运行环境的版本表对应这条生成链，便于沿产物关系定位问题。\par


目标程序采用另一条工具链。RISC-V GCC编译程序与所需运行支持，链接器确定代码和数据布局，binutils提供二进制转换及反汇编。NEMU reference在宿主侧构建为可加载的参考库，由EMU运行参数指定。这样，硬件仿真器、目标程序和参考模型具有各自清楚的构建入口，并在差分运行时汇合。\par


实际组织时，构建日志与运行日志分别保存。构建记录说明工具如何生成产物；运行记录说明加载器选择了哪些文件、程序如何启动以及最后如何结束。两者通过产物路径、文件摘要和运行参数对应。对一个具体问题，先找到相关产物，再沿依赖链向上查看生成步骤，能够较快区分源码、工具和装载配置三个层面的原因。\par


{\small\setlength{\tabcolsep}{3pt}\begin{longtable}{>{\raggedright\arraybackslash}p{0.2\linewidth}>{\raggedright\arraybackslash}p{0.24\linewidth}>{\raggedright\arraybackslash}p{0.24\linewidth}>{\raggedright\arraybackslash}p{0.22\linewidth}}
\caption{构建链路与运行产物}\label{tab:1-1}\\
\toprule
\textbf{产物} & \textbf{生成入口} & \textbf{运行中的作用} & \textbf{主要定位信息}\\\midrule\endfirsthead
\caption[]{构建链路与运行产物（续）}\\
\toprule
\textbf{产物} & \textbf{生成入口} & \textbf{运行中的作用} & \textbf{主要定位信息}\\\midrule\endhead
{香山EMU} & {硬件生成与宿主编译} & {执行DUT仿真} & {硬件版本、配置、构建参数、文件摘要}\\
{NEMU reference} & {参考模型构建} & {提供架构状态参考} & {源码版本、配置、加载路径、文件摘要}\\
{程序ELF/BIN} & {RISC-V编译与链接} & {执行hello及vadd测试} & {入口、装载段、反汇编、镜像摘要}\\
{共同flash} & {启动汇编与二进制转换} & {建立一致的启动输入} & {文件长度、内容摘要、显式加载参数}\\
\bottomrule\end{longtable}}

\reportfigure{S1-01}{开发环境与软硬件协同验证框架}{12.1cm}{Windows宿主中的WSL2环境承载构建与差分运行。实线为构建依赖，虚线为运行使用；程序BIN与启动flash是共同输入，EMU和NEMU参考库分别提供DUT与REF功能。框图不表示独立操作系统进程或内部逐周期锁步。阶段一历史环境，未混入最终0f身份。}{精确SVG工程图；工程依据：环境记录、\seqsplit{team/scripts/stage1.sh}及已固定版本表；原图存档与所选资产身份见\seqsplit{figures/redesigned-20261105/selected-diagrams.json}。阶段一历史环境；与最终0f组合分别标注}

\section{环境部署与运行步骤}

\subsection{WSL资源配置}

本次为 WSL 分配 20 GB 内存和 16 GB 交换空间。Windows 用户目录 \%UserProfile\% 下的 .wslconfig 内容如下：\par


\Needspace{5\baselineskip}
\begin{Verbatim}[breaklines,breakanywhere,breaksymbolleft={},breaksymbolright={},fontsize=\footnotesize,baselinestretch=1]
[wsl2]
memory=20GB
swap=16GB
\end{Verbatim}


保存正在进行的 WSL 工作后，在 PowerShell 执行 wsl -{}-shutdown，重新进入 Ubuntu，再用 free -h 核对生效结果。原机器的交换文件路径为 D:\textbackslash{}WSL\textbackslash{}wsl-swap.vhdx；新环境按本机磁盘布局选择路径即可。\par


内存规划需要同时考虑 JVM、硬件生成工具和并行 C++ 编译进程。交换空间可为短时内存压力提供缓冲，但频繁换页会增加构建时间，因此它与可用物理内存具有不同作用。内存较小的机器应相应降低 JVM 堆与编译并行度，使多个进程的合计需求处于可用范围内。构建日志同时记录并行参数和结束状态，便于区分资源配置问题与源码编译问题。\par

\subsection{基础工具安装}

完成配套源码检出后，在团队仓库根目录加载路径，并安装基础工具：\par


\Needspace{7\baselineskip}
\begin{Verbatim}[breaklines,breakanywhere,breaksymbolleft={},breaksymbolright={},fontsize=\footnotesize,baselinestretch=1]
source team/scripts/activate.sh
sudo bash "$TEAM_XS_ENV_ROOT/setup-tools.sh" --target base
sudo bash "$TEAM_XS_ENV_ROOT/setup-tools.sh" --target llvm
sudo bash "$TEAM_XS_ENV_ROOT/setup-tools.sh" --target mill
sudo apt-get install -y cmake ninja-build gtkwave
\end{Verbatim}


入口来自固定 xs-env 的 \href{https://github.com/OpenXiangShan/xs-env/blob/b0b506481e839a5df8a580e76f5aa246dc60fe7c/setup-tools.sh}{setup-tools.sh}。软件源提供的补丁版本可能随时间变化，安装步骤完成后仍需通过环境检查确认实际版本。这样可以把“安装入口相同”与“得到的工具版本相同”分别记录。\par

\subsection{JDK安装与校验}

原安装入口使用 Oracle 21/latest 地址。为了恢复本阶段的 21.0.12，团队固定 GraalVM 安装脚本使用对应 Linux x64 归档及官方 SHA-256，安装到用户目录：\par


\Needspace{5\baselineskip}
\begin{Verbatim}[breaklines,breakanywhere,breaksymbolleft={},breaksymbolright={},fontsize=\footnotesize,baselinestretch=1]
export JAVA_HOME="$HOME/.local/share/graalvm-jdk-21.0.12"
bash team/scripts/install-graalvm.sh
source team/scripts/activate.sh
\end{Verbatim}


归档校验用于确认下载字节，版本检查用于确认实际启动的 Java。二者配合，可以发现下载路径正确但环境变量仍指向另一套 JDK 的情况。新终端继续导出相同 JAVA\_HOME。原环境中已有 \seqsplit{/opt/graalvm-jdk-21} 的安装，也可以在确认版本后沿用；脚本遇到版本不同的已有安装会停止，以保留当前配置。\par


安装方法参考 \href{https://docs.oracle.com/en/graalvm/jdk/21/docs/getting-started/installation-linux/}{Oracle GraalVM 安装说明}。固定下载项为 \href{https://download.oracle.com/graalvm/21/archive/graalvm-jdk-21.0.12\_linux-x64\_bin.tar.gz}{21.0.12 Linux x64}，校验来源为\href{https://download.oracle.com/graalvm/21/archive/graalvm-jdk-21.0.12\_linux-x64\_bin.tar.gz.sha256}{官方 SHA-256}。工具更新时同步记录新版本及相应验证结果，便于将新产物与本阶段历史产物区分。\par

\subsection{Verilator安装与环境核查}

Verilator 从固定提交构建，以下流程显式使用两个编译任务。构建并行度应与本机可用资源相配合；阶段一 EMU 编译记录中的四任务参数与这里安装工具本身的参数分别记录。\par


\begin{Verbatim}[breaklines,breakanywhere,breaksymbolleft={},breaksymbolright={},fontsize=\footnotesize,baselinestretch=1]
sudo apt-get install -y help2man perl autoconf flex bison libgoogle-perftools-dev \
  libjemalloc-dev numactl perl-doc zlib1g-dev liblz4-dev libfl2 libfl-dev
mkdir -p "$HOME/riscv-tools"
git clone https://github.com/verilator/verilator.git "$HOME/riscv-tools/verilator"
cd "$HOME/riscv-tools/verilator"
git checkout ea338be98e1e838d3518809ce8899f85a009963c
unset VERILATOR_ROOT
autoconf
./configure CC=/usr/local/llvm/bin/clang CXX=/usr/local/llvm/bin/clang++ LINK=/usr/local/llvm/bin/clang++
make -j2
sudo make install
\end{Verbatim}


编译前清除旧的 VERILATOR\_ROOT，并显式选择已安装的 Clang，以便构建输入与记录一致。安装后回到团队仓库根目录检查环境：\par


\Needspace{4\baselineskip}
\begin{Verbatim}[breaklines,breakanywhere,breaksymbolleft={},breaksymbolright={},fontsize=\footnotesize,baselinestretch=1]
source team/scripts/activate.sh
bash team/scripts/check-environment.sh
\end{Verbatim}


环境检查核对配套源码提交、香山递归子模块和关键工具版本。香山提交可以包含本队新增文档，但应保留指定硬件基线祖先。检查结果用于确认构建条件；硬件行为则由对应版本的仿真记录说明。这样的分工使环境恢复与功能验证各自具有清楚的输出，也便于定位失败发生在哪一层。\par

\subsection{程序构建环境}

原记录的工作目录为 \seqsplit{/home/riscv/riscv-lab/xs-env/XiangShan}。复用其他机器或其他检出时，在实际团队仓库根目录加载环境即可，并保留本次路径和提交记录。\par


activate.sh 将当前团队仓库设置为 NOOP\_HOME，按照 \seqsplit{TEAM\_XS\_ENV\_ROOT} 指向配套 NEMU 和 AM，并设置 \seqsplit{LINUX\_GNU\_TOOLCHAIN=1}。这些变量把硬件仓库、参考模型和程序构建入口连接起来，运行前可据此核对各组件来自哪套环境。\par


AM 构建向量工作负载时仍需显式指定 MARCH=rv64gcv。这使编译器采用包含向量扩展的目标 ISA，与 hello-vadd 所用的向量装载、加法和存储指令一致。程序构建完成后，将 ELF、运行镜像、反汇编和实际运行命令一起归档，便于从源程序追踪到日志中的提交 PC。\par

\subsection{集成构建与运行入口}

原实验使用Windows下的WSL2 Ubuntu 22.04.5，宿主处理器为Ryzen 7 7840H，可用物理内存约27.7GiB。WSL配置20GB内存与16GB swap，硬件生成、编译和仿真在WSL原生文件系统中执行。下表列出阶段一的主要版本与运行设置；完整子模块锁和工具输出列于文后。\par

{\small\setlength{\tabcolsep}{3pt}\begin{longtable}{>{\raggedright\arraybackslash}p{0.25\linewidth}>{\raggedright\arraybackslash}p{0.69\linewidth}}
\caption{集成构建与运行入口}\label{tab:1-2}\\
\toprule
\textbf{组件} & \textbf{固定版本或设置}\\ \midrule\endfirsthead
\caption[]{集成构建与运行入口（续）}\\
\toprule
\textbf{组件} & \textbf{固定版本或设置}\\ \midrule\endhead
xs-env & \seqsplit{b0b506481e839a5df8a580e76f5aa246dc60fe7c}\\
原香山基线 & \seqsplit{16ae9ddcda54fc9a2fddffad73174cf793ac7814}\\
原NEMU基线 & \seqsplit{a48a11f208ff81ad765591d83271d1b339fb3861}\\
nexus-am & \seqsplit{086a238062fdfb3281b745a537b361983d0154b5}\\
JDK / Mill / firtool & 21.0.12 / 0.12.3 / 1.62.1\\
Verilator & 5.052，固定ea338be提交\\
RISC-V GCC / binutils & 11.4.0 / 2.38\\
DUT & MinimalConfig、单核、VLEN128\\
原波形构建 & EMU\_THREADS=1、EMU\_TRACE=fst\\
\bottomrule\end{longtable}}

部署按照依赖关系推进：配置WSL资源，检出配套源码和递归依赖，核对工具版本，分别构建NEMU reference与香山EMU，编译hello/vadd及启动镜像，最后执行差分运行。命令、退出状态、版本和文件摘要随各阶段产物一起保存。\par

复现入口由以下脚本组成。工作负载中的团队标识由成员确认后填写；构建与运行输出应使用独立目录，以便对应具体实验。\par

\begin{Verbatim}[breaklines,breakanywhere,breaksymbolleft={},breaksymbolright={},fontsize=\footnotesize,baselinestretch=1]
source team/scripts/activate.sh
bash team/scripts/check-environment.sh
bash team/scripts/stage1.sh nemu
bash team/scripts/stage1.sh emu
TEAM_ID=正式团队标识 bash team/scripts/stage1.sh workload
bash team/scripts/stage1.sh run
\end{Verbatim}

部署中遇到的首个差分问题位于启动PC 0x10000004。检查发现两端默认flash分别包含3条和6条启动指令。随后使用同一boot-flash.S生成12字节镜像，并在启动命令中显式指定共同flash、128MB RAM和固定reference。完整定位过程见后文“启动差分故障的定位与修复”。\par

\reportfigure{S1-02}{环境部署与验证执行流程}{14.9cm}{环境准备、工具核查、产物构建、DiffTest运行和成功归档构成主流程；失败证据另行归档。定位修正后按故障原因返回准备、工具核查或构建环节，虚线为处理路径，不表示每一类历史故障均实际发生，也不表示脚本自动诊断或修复。2026-10-09阶段一流程：香山16ae9dd、NEMU a48a11f。}{精确SVG工程图；工程依据：阶段一真实执行顺序、构建和运行日志；原图存档与所选资产身份见\seqsplit{figures/redesigned-20261105/selected-diagrams.json}。2026-10-09阶段一流程}

\subsection{构建资源与并行恢复}

香山EMU的构建涉及多个资源特征不同的阶段。Java进程需要堆空间解析工程和生成硬件，firtool处理硬件表示，Verilator生成仿真代码，C++阶段则启动多个编译进程并产生对象文件。宿主文件缓存与操作系统也占用内存。资源配置要为这些部分保留共同工作的空间，单独增大某一进程的内存上限并不等于增加了整条构建链的可用资源。\par


原实验主机可用物理内存约27.7GiB，WSL配置20GB内存和16GB swap。这里的20GB是WSL内存上限，swap用于承接部分内存压力；二者在访问成本和作用上不同。持续发生换页时，构建墙钟时间可能明显增加。因此，资源记录同时关注内存配置、并行度、进程状态和日志进展，而不是只记录一项内存容量。\par


2026年10月9日的构建记录中，C++编译在20:42:15中断并恢复，并行度调整为4，已完成对象保留；EMU随后在20:53:32完成构建。这个过程体现了分阶段调节并发的方式。复现脚本中的较低默认并发适合控制初次构建的内存压力，原记录中的-j4则对应当时的实际执行参数。两种配置都通过各自命令和时间记录定位。\par


并行资源也应与仿真配置分开描述。阶段一DUT采用单核，原波形构建使用单仿真线程；宿主可以同时运行多少独立任务，属于另一项调度条件。构建并发影响宿主资源和等待时间，程序的EMU周期统计则来自仿真内部。将构建耗时、宿主执行时间和DUT周期分别记账，便于后续评价环境效率与指令性能。\par


对象文件和构建缓存有助于恢复长作业，但它们对应具体源码、工具和编译参数。沿用已完成对象时，继续保存原命令与依赖身份；调整会影响生成内容的参数时，按照构建系统的依赖关系更新相应产物。报告保留中断和恢复记录，使最终EMU的形成过程具有连续的解释，而非只有最后一条成功提示。\par

\begin{itemize}
\item 2026年10月9日20:42:15，中断并恢复C++构建，将并行度调为4，保留已完成对象。
\item 2026年10月9日20:53:32，EMU构建结束，状态文件记录exit\_code=0。
\item 2026年10月9日20:55:26，workload与共同flash构建结束，记录\seqsplit{WORKLOAD\_BUILD\_PASS}。
\item 原正常运行于21:00:28结束，退出状态为0。
\end{itemize}

\subsection{程序镜像、装载与启动}

ELF、BIN和flash在运行中承担不同作用。ELF包含符号、代码与数据段信息，可用于查看入口地址和反汇编；BIN提供加载器使用的程序字节；flash提供复位后的启动代码。启动镜像与目标程序是相互配合的两个输入，分析程序尚未到达C入口的差异时，flash及装载配置通常是重要线索。\par


反汇编连接源程序和硬件观察点。编译器选择指令，链接器安排代码与数据，最终地址由这组构建条件共同决定。通过目标指令的PC查找同次构建的机器字，再与提交记录对应，可以确定波形中观察的是哪一次操作。即使C代码保持原样，链接参数或依赖库改变后，PC表也可能随之变化。\par


显式运行参数把构建产物与一次执行绑定。阶段一在命令中指定共同flash、128MB RAM与固定reference，使两端启动输入可以直接检查。共同文件的含义是字节实物一致；两端如何将它装入各自平台地址空间，仍由各自装载实现决定。描述共同启动条件时，文件内容、装载参数和地址映射分别列出，因果关系更清楚。\par


hello程序提供了启动链路中的一个软件观察点。串口输出表明程序已进入对应的软件路径，随后vadd用例设置向量状态、执行运算并打印比较结果。原hello中的IPv4来自WSL宿主，用于识别当时环境；团队字符串为原记录中的riscv-vdot-team。正式身份信息与历史输出分别保存，便于在更新展示材料时继续检索原运行。\par


装载问题与运算问题也具有不同的日志特征。启动日志给出flash、程序和reference的选择，软件输出给出已到达的程序阶段，差分信息给出首次观察到的架构差异。按执行顺序阅读这些信息，可以先确定故障发生在启动之前、启动过程中还是进入测试程序之后，再选择相应源码和数据继续分析。\par

\section{启动差分故障的定位与修复}

原失败的首次差异位于PC 0x10000004，涉及t0寄存器，记录的instrCnt为2。这个位置处于启动阶段，早于hello/vadd测试主体。定位首先围绕两端的初始执行序列展开：检查启动输入，查看该位置的机器字，并与启动文件的反汇编对应。\par


检查发现，两端默认flash分别包含3条和6条启动指令。默认输入差异使两端在到达同一程序入口前就经历了不同的执行过程。修复采用一份boot-flash.S生成共同的12字节boot-flash.bin，再通过显式flash参数将同一实物交给DUT与NEMU。三条启动指令的文件偏移、机器字与反汇编保留在图数据附件中。\par


文件偏移定位flash中的指令，启动PC定位处理器执行地址，两者通过装载映射关联。同一份12字节boot-flash.bin分别交给DUT与NEMU装载。首次差异PC用于定位原故障，后续终止或ABORT事件的地址按对应事件记录。\par


恢复运行继续启用DiffTest，执行三类vadd用例，记录24个元素通过、GOOD TRAP和exit 0。失败记录与恢复记录分别保存，前者解释问题如何被发现，后者说明统一启动输入之后的结果。报告中的修复过程由“首次差异—默认输入检查—共同文件—正常运行”组成，每一步都可以回到原日志或构建实物。\par


{\small\setlength{\tabcolsep}{3pt}\begin{longtable}{>{\raggedright\arraybackslash}p{0.27\linewidth}>{\raggedright\arraybackslash}p{0.3\linewidth}>{\raggedright\arraybackslash}p{0.35\linewidth}}
\caption{启动差分故障的定位与修复}\label{tab:1-3}\\
\toprule
\textbf{排查步骤} & \textbf{使用材料} & \textbf{定位结果}\\\midrule\endfirsthead
\caption[]{启动差分故障的定位与修复（续）}\\
\toprule
\textbf{排查步骤} & \textbf{使用材料} & \textbf{定位结果}\\\midrule\endhead
{查找首次差异} & {原失败日志、PC与t0状态} & {差异发生于启动阶段}\\
{比较启动输入} & {默认flash与反汇编} & {两端默认指令序列不同}\\
{统一启动文件} & {同一boot-flash.bin及显式参数} & {两端加载同一12字节实物}\\
{查看恢复结果} & {软件摘要、GOOD TRAP与退出状态} & {hello/vadd完成正常差分运行}\\
\bottomrule\end{longtable}}


这一案例也给后续故障分析提供了合适的起点：先定位最早出现差异的阶段，再分析对应输入和执行状态。若差异仍在启动入口，应优先检查flash与装载；若已经进入向量程序，则转向指令、寄存器和结果记录。分层定位可以缩小检查范围，并使每次参数调整与观察结果直接对应。\par

\reportfigure{S1-07}{启动输入统一与差分恢复}{12.1cm}{首次差异位于启动PC 0x10000004，instrCnt为2。共同boot-flash.bin的12字节由三条机器指令组成，图中0–3、4–7、8–11为文件字节偏移。显式向RTL与NEMU提供同一启动文件后，三类共24个元素通过检查，以GOOD TRAP和exit 0结束；共享文件不等于两端物理映射地址相同。}{2026-10-09原失败日志行164/166、\seqsplit{boot-flash-disassembly.txt}、完整成功日志行7/10/16/17/19及\seqsplit{hello-vadd-summary.log}；完整出处和哈希见图数据附件。历史启动问题定位；不是vdot错误或新运行。共用文件不等于两端物理映射地址相同}

\section{hello与vadd运行结果}

hello输出包含团队标识和WSL宿主IPv4。2026年10月9日记录使用的团队字符串为riscv-vdot-team，本文按原日志保留。正式团队名称等身份信息集中列于封面。\par

vadd程序采用SEW8、LMUL1和VL8，依次执行两次向量加载、一次vadd.vv及一次向量存储，再将八个结果字节与标量参考逐项比较。三个用例分别检查基本加法、负数补码和定宽回绕。\par

{\small\setlength{\tabcolsep}{3pt}\begin{longtable}{>{\raggedright\arraybackslash}p{0.25\linewidth}>{\raggedright\arraybackslash}p{0.69\linewidth}}
\caption{hello与vadd运行结果}\label{tab:1-4}\\
\toprule
\textbf{用例} & \textbf{输出字节}\\ \midrule\endfirsthead
\caption[]{hello与vadd运行结果（续）}\\
\toprule
\textbf{用例} & \textbf{输出字节}\\ \midrule\endhead
正数 & 9 9 9 9 9 9 9 9\\
负数补码 & 254 252 250 248 246 244 242 240\\
边界回绕 & 128 127 0 0 0 0 254 254\\
\bottomrule\end{longtable}}

正常运行记录包含HELLO\_VADD\_PASS cases=3 lanes=24、GOOD TRAP和真实退出状态0，运行期间启用DiffTest。EMU统计为16557条指令、42962周期，统计对象是包含启动、UART与结果校验的完整程序。各用例的数值解释见本节“定宽加法与字节结果”，事件关系见“基于波形的向量执行过程”。\par

\reportfigure{S1-03}{hello与vadd正常结束输出}{8.5cm}{运行日志记录三类vadd输出、DiffTest比较和GOOD TRAP终态。}{用户截图 \seqsplit{WindowsTerminal\_m5hXeDJnV1.png}；原\seqsplit{hello-vadd-summary.log}。2026-10-09历史运行回放，不是本次重新运行}

\subsection{10月25日英文输出与运行验证}

2026年10月25日（UTC+8），补证实验使用个人标识LiZufan001构建hello/vadd程序，身份状态为configured。原始串口输出如下，其中IPv4为构建宿主地址：\par
\begin{Verbatim}[breaklines,breakanywhere,breaksymbolleft={},breaksymbolright={},fontsize=\footnotesize,baselinestretch=1]
hello xiangshan, I am LiZufan001, IP address=172.27.184.54
identity_status=configured; IP address is build-host IPv4
\end{Verbatim}
同一BIN分别在已封存生产EMU和独立FST模型上启用普通NEMU DiffTest，两次运行均通过3组、24个元素检查，记录\Verb|HELLO_VADD_PASS cases=3 lanes=24|、GOOD TRAP及退出状态0。每次统计均为16621条指令、43249周期，覆盖启动、UART和结果校验的整程序。\href{https://github.com/cc77-builder/riscv/blob/2690c7370fb7c94111b4619bbac352e9be4ccc4d/team/evidence/stage1-current-20261025/production-run/emulator-output.log}{原始输出}与\href{https://github.com/cc77-builder/riscv/blob/2690c7370fb7c94111b4619bbac352e9be4ccc4d/team/evidence/stage1-current-20261025/production-run/run.json}{运行回执}保存了具体输入、命令和结束状态。\par

\subsection{定宽加法与字节结果}

阶段一vadd采用SEW8、LMUL1和VL8，参与用例比较的是八个8位元素。程序通过两次向量加载取得输入，执行按元素加法，再将结果存回内存供标量比较。三个用例使用相同的执行入口，分别选择便于观察基本运算、补码和回绕行为的数据，从而让环境验证具有不同的数值特征。\par


\href{https://github.com/riscvarchive/riscv-v-spec/blob/v1.0/v-spec.adoc}{RISC-V V1.0规范}定义同宽向量整数加法，并说明未特别指定的整数运算在溢出时回绕。八位加法因此按模256保存结果。若把一个结果字节记为r，它的无符号表示位于0到255之间；按八位补码解释时，最高位为1的字节对应负值。例如原负数用例打印的254是十六进制FE，对应补码整数−2。打印形式由格式选择决定，存储的八位比特保持相同。比较时以相同宽度和截断规则计算标量参考，就能直接核对程序结果。\par


正数用例的八个结果字节均为9，低64位显示为\seqsplit{0x0909090909090909}。负数用例按元素顺序打印254、252、250、248、246、244、242、240，对应低64位\seqsplit{0xf0f2f4f6f8fafcfe}。观察十六进制向量值时，右端是最低字节，因而与按元素序号递增打印的字节列表方向不同。把两个显示形式按字节位置对应，可以避免把书写顺序差异当成数据差异。\par


边界用例输出128、127、0、0、0、0、254、254，低64位为\seqsplit{0xfefe000000007f80}。这里同时出现最高位变化、零和接近字节上限的值，适合检查定宽结果的表示。八个字节在向量低64位中的组合关系可以写成各结果字节乘以相应的256次幂后求和；因此逐元素列表与整体十六进制值是同一结果的两种表达。\par


VLEN128说明向量寄存器的物理宽度，VL8说明该用例的活动元素数。本文结果表按实际八个元素解释，整体寄存器宽度与有效运算范围分别标注。随后vdot使用低八字节参与点积并形成自定义结果布局，属于另一套指令语义，其低32位结果和高位规则在阶段二说明。\par


软件输出还保留了用例数量与比较数量。HELLO\_VADD\_PASS中的cases=3对应三种输入组织，lanes=24对应三次各八个活动元素。读者可以将该摘要与逐字节结果及GOOD TRAP、exit 0一起阅读，得到从数值比较到运行结束的完整解释。\par

\section{基于波形的向量执行过程}

取指和译码之后，向量指令经重命名、调度与操作数准备进入执行单元。两个vle8.v提供输入，vadd.vv完成按元素加法，结果进入写回路径并在ROB提交。差分系统在架构状态比较点与NEMU对应状态核对，内部观察器则记录输入、输出和提交事件。\par

原波形提取的vadd提交PC为0x80000138，机器码为0x021101d7，vecwen=1、wdest=3。下表按这组指令身份关联三个用例的事件坐标与低64位结果。\par

{\small\setlength{\tabcolsep}{3pt}\begin{longtable}{>{\raggedright\arraybackslash}p{0.1\linewidth}>{\raggedright\arraybackslash}p{0.16\linewidth}>{\raggedright\arraybackslash}p{0.16\linewidth}>{\raggedright\arraybackslash}p{0.16\linewidth}>{\raggedright\arraybackslash}p{0.34\linewidth}}
\caption{基于波形的向量执行过程}\label{tab:1-5}\\
\toprule
\textbf{用例} & \textbf{输入事件坐标} & \textbf{输出事件坐标} & \textbf{提交事件坐标} & \textbf{运算输出低64位}\\ \midrule\endfirsthead
\caption[]{基于波形的向量执行过程（续）}\\
\toprule
\textbf{用例} & \textbf{输入事件坐标} & \textbf{输出事件坐标} & \textbf{提交事件坐标} & \textbf{运算输出低64位}\\ \midrule\endhead
正数 & 19625 & 19626 & 19631 & \seqsplit{0x0909090909090909}\\
负数 & 25460 & 25461 & 25466 & \seqsplit{0xf0f2f4f6f8fafcfe}\\
边界 & 33009 & 33010 & 33015 & \seqsplit{0xfefe000000007f80}\\
\bottomrule\end{longtable}}

三个观察点展示了执行结果产生到架构提交的流水关系。表中数字是原CSV事件坐标，图中间距用于组织阅读；现有局部截图未显示时钟，因而本节按事件位置解释记录，不把坐标差换算为物理时间或固定指令时延。需要进一步查看握手时，可恢复对应完整波形，联合定位时钟、valid、ready和结果信号。\par

{\raggedright 原归档记载的完整FST路径为\seqsplit{/home/riscv/riscv-lab/stage1-evidence/hello-vadd.fst}，文件大小220295598字节。该记录属于2026年10月9日基线，恢复时以归档索引和文件摘要定位。最终0f环境的hello/vadd验证记录由R5验收单独绑定。\par}

\reportfigure{S1-06}{vadd输入、输出与提交事件}{8.2cm}{三类vadd用例通过PC 0x80000138、指令0x021101D7和目的v3关联。数字为CSV事件坐标。}{\seqsplit{evidence/vector-fix-alu-events.csv}及\seqsplit{vadd-commit-events.csv}；生成源与哈希见\seqsplit{figures/evidence-20261105}。2026-10-09历史vadd事件；不能推导物理时钟或新硬件延迟}

\subsection{执行、写回与提交}

程序中的一条向量指令会经过多个硬件阶段。取指和译码确定指令及操作类型，重命名与调度安排数据依赖，执行单元计算结果，写回路径传递结果与身份，ROB按架构顺序提交。执行结果出现与架构提交是不同观察点，二者之间还存在父级选择及在途状态的处理。\par


阶段一按PC 0x80000138、机器字0x021101d7及目的v3关联vadd记录。提交侧vecwen=1说明该记录带有向量写入，wdest=3给出对应目的。PC和机器字共同提供静态指令位置，目的信息帮助确认被观察的结果。它们与用例顺序、事件坐标和结果值结合后，才能把三类软件输入准确连接到各自的硬件事件。\par


正数用例的输入、输出和提交事件分别位于19625、19626和19631，负数用例为25460、25461和25466，边界用例为33009、33010和33015。这些数字沿用原CSV坐标，表示所列记录在波形中的位置。图中采用三列布局突出阶段关系，列间距离没有按原坐标比例绘制。\par


内部信号的含义还取决于采样位置。输入记录可能表示执行端可见的操作数，输出记录表示结果观察点，握手则需要结合valid和ready判断一次事务是否实际转移。阶段一图表按现有输入、输出与提交字段命名；需要讨论握手行为时，应在原波形对应位置同时查看双方条件。这样既能利用已有记录解释流水关系，也能清楚说明各字段承担的作用。\par


现有局部波形截图未显示时钟。截图中的ps是波形显示单位，事件坐标是用于定位的数值；EMU周期统计和宿主运行时间则来自不同的计量位置。本阶段通过PC、字段与坐标连接记录，没有为这些局部窗口建立物理频率或固定时延结论。\par


乱序执行中，程序依赖、执行资源及下游状态会影响执行和提交的间隔。报告用三个具体窗口说明该次运行的过程，后续性能分析则采用专门的DUT计数窗口和相同工作量对照。把事件定位与性能计量分开，读者就能分别回答“结果在哪里出现”和“规定工作用了多少周期”。\par

\reportfigure{S1-04}{vadd正数用例的执行与提交波形}{16.0cm}{(a)输出位置及相关事件；(b)输入游标处信号；(c)提交游标处信号。19625、19626和19631为selected VCD中的事件坐标，截图单位为ps，时钟信号未显示。}{用户上传正数输出截图及\seqsplit{riscv-lab/picture}中的输入、提交截图；保留原始图并仅裁剪外部空白。2026-10-09原基线vadd；此局部截图未显示时钟，2026-10-25 T04含时钟补证另见正文；不是0f vdot}

\reportfigure{S1-05}{vadd负数与边界用例结果波形}{10.0cm}{负数与边界样例的输出低64位分别为\seqsplit{F0F2F4F6F8FAFCFE}和\seqsplit{FEFE000000007F80}，游标显示各自输入、输出与提交位置。}{用户上传 \seqsplit{msrdc\_yWBpPeYM8q.png} 与 \seqsplit{msrdc\_6cIzJ9aIVJ.png}。同一2026-10-09历史vadd波形；低64位按有效8个元素解释}

\subsection{指令、事件与波形联合定位}

波形、CSV与反汇编提供互补的观察方式。反汇编给出静态地址与机器字，事件CSV按字段记录被提取的变化，FST保留更完整的信号历史。局部截图适合快速呈现某一窗口，完整文件则适合围绕具体疑点展开查看。报告按这些用途组织图、表与附件。\par


一种直接的定位顺序是先从提交记录取得PC，再在对应反汇编中找到指令，随后用事件坐标打开执行窗口。对vadd，还可以核对机器字、目的v3和结果低64位；如果三者与软件用例一致，就可以进一步沿信号层级查看操作数和控制条件。这条路径把静态程序、软件输出和内部波形连接到同一条被观察指令。\par


阅读结果值时需要注意数据展示方向。日志可以按元素顺序输出十进制字节，波形窗口也可以把一组字节合并显示为十六进制。通过元素宽度和低位位置统一表示后，正数、负数和边界三个用例均可以按同样步骤阅读。数据解释与时间定位在这里是两项独立工作：先确认数值对应，再查看它出现在哪个观察点。\par


信号索引记录了提交与执行端的信号层级、位宽和标识，帮助读者解释CSV中的字段。检查某个字段时，可结合索引确定原信号；恢复对应完整波形后，再查看该坐标附近的上下文。报告通过精选窗口与附件索引兼顾主要过程的阅读和具体疑点的深入分析。\par

\subsection{10月25日含时钟的波形调试记录}

10月25日的独立FST模型使用同一组既有硬件生成物构建，保留真实时钟信号。最终提取覆盖两个VialuFix、八个提交槽和127个信号，检查33249个完整上升沿，关联三次FU输入、三次FU输出与三次DPIC提交观测，共9条事件。\par
采样取每个局部上升沿前的稳定低相值，FST原始时间与周期计数分别保存。DPIC记录经过Rob.scala中delay=3的观察路径，其坐标表示提交接口的观测时刻。选通信号FST可在GTKWave中查看时钟、操作数、握手与事务标识。本节原波形图仍对应10月9日；10月25日记录由独立\href{https://github.com/cc77-builder/riscv/blob/2690c7370fb7c94111b4619bbac352e9be4ccc4d/team/evidence/stage1-current-20261025/wave/wave.json}{波形与事件索引}定位。\par
该次构建的向量ALU接口见固定源码中的\href{https://github.com/cc77-builder/riscv/blob/34bae6d4acceac939f958199dae53c14bc1f6183/src/main/scala/xiangshan/backend/fu/wrapper/VIAluFix.scala}{VIAluFix包装层}，用于对照操作数、控制与结果的连接。\par

\section{阶段结论与复现资料}

\subsection{阶段结论与材料入口}

阶段一完成了配套环境部署、启动差异定位、共同flash配置和三类向量加法用例的差分运行。软件输出、EMU退出状态、事件CSV和原波形为后续分析提供了互相对应的记录；自定义vdot的语义、父级控制与性能验证在后续阶段展开。\par

阶段一证据目录为\seqsplit{team/evidence/stage1-20261009}，包含\seqsplit{hello-vadd-summary.log}、工具版本、子模块锁、boot-flash反汇编、\seqsplit{vadd-commit-events.csv}及\seqsplit{vector-fix-alu-events.csv}。程序源码位于\seqsplit{team/workloads/hello-vadd}，本文源目录中的evidence保存用于报告展示的对应材料。\par

\href{https://github.com/cc77-builder/riscv/blob/eca41d7cd959f5f3c4f240e382919c492d1ccb96/team/docs/stage1.md}{阶段一原记录} · \href{https://github.com/cc77-builder/riscv/blob/eca41d7cd959f5f3c4f240e382919c492d1ccb96/team/docs/environment.md}{环境记录}\par

\Needspace{8\baselineskip}
\subsection{源码版本与依赖恢复}

本阶段使用的配套源码由以下完整提交号标识。xs-env 提供环境入口，香山及其子模块决定硬件生成输入；NEMU、nexus-am 和 DRAMsim3 分别参与参考执行、程序构建和存储系统仿真。对这些组件同时记录版本，可以避免仅固定顶层仓库、依赖却随下载时间变化的问题。\par


{\small\setlength{\tabcolsep}{3pt}\begin{longtable}{>{\raggedright\arraybackslash}p{0.29\linewidth}>{\raggedright\arraybackslash}p{0.65\linewidth}}
\caption{源码版本与依赖恢复}\label{tab:1-6}\\
\toprule
\textbf{组件} & \textbf{已验证提交}\\\midrule\endfirsthead
\caption[]{源码版本与依赖恢复（续）}\\
\toprule
\textbf{组件} & \textbf{已验证提交}\\\midrule\endhead
{xs-env} & {\seqsplit{b0b506481e839a5df8a580e76f5aa246dc60fe7c}}\\
{香山硬件基线} & {\seqsplit{16ae9ddcda54fc9a2fddffad73174cf793ac7814}}\\
{NEMU} & {\seqsplit{a48a11f208ff81ad765591d83271d1b339fb3861}}\\
{nexus-am} & {\seqsplit{086a238062fdfb3281b745a537b361983d0154b5}}\\
{DRAMsim3} & {\seqsplit{1d2a9bf6da8e050d975ef2669c42faa93943f489}}\\
{DiffTest} & {\seqsplit{62e4540c03c03aa947bdb67f4b6bf66f1a29c5d0}}\\
{yunsuan} & {\seqsplit{cadd3c2f43096e253d61296b33ac697be8354e29}}\\
\bottomrule\end{longtable}}


在团队仓库执行 bash \seqsplit{team/scripts/prepare-dependencies.sh}。该入口检出固定 xs-env，初始化配套 NEMU、nexus-am、DRAMsim3 及其递归子模块，并锁定 NEMU 的额外依赖。团队仓库自身的香山子模块通过 git clone \seqsplit{-{}-recurse-submodules} 获取。脚本在发现已有依赖版本不一致或已跟踪文件有修改时停止，保留现场供核对，避免把不同输入混入一次恢复过程。\par


NEMU 额外依赖的提交记录见锁定脚本。其中 SimPoint 按 NEMU 已记录的子模块版本递归初始化，其余下载仓库也以锁定提交为准。初始化完成后，应核对实际检出的提交，而不是只确认下载命令返回成功；完整源码锁与后续构建日志共同说明产物的来源。\par

\subsection{完整工具版本}

工具链包括运行构建系统的 JDK、硬件生成与转换工具、宿主编译器、目标程序编译器，以及波形阅读工具。各类工具处理的输入不同，记录版本时也应保留完整输出；例如目标 RISC-V GCC 决定程序指令，而宿主 Clang 或 GCC 处理的是仿真器的 C++ 代码。\par


{\small\setlength{\tabcolsep}{3pt}\begin{longtable}{>{\raggedright\arraybackslash}p{0.29\linewidth}>{\raggedright\arraybackslash}p{0.65\linewidth}}
\caption{完整工具版本}\label{tab:1-7}\\
\toprule
\textbf{工具} & \textbf{版本}\\\midrule\endfirsthead
\caption[]{完整工具版本（续）}\\
\toprule
\textbf{工具} & \textbf{版本}\\\midrule\endhead
{Oracle GraalVM JDK / javac} & {21.0.12}\\
{LLVM / Clang} & {19.1.7}\\
{Mill 项目版本} & {0.12.3}\\
{Scala / Chisel} & {2.13.15 / 6.6.0，由项目解析}\\
{firtool} & {1.62.1，由本次项目依赖解析}\\
{Verilator} & {5.052，提交 \seqsplit{ea338be98e1e838d3518809ce8899f85a009963c}}\\
{RISC-V GCC} & {11.4.0，\seqsplit{riscv64-linux-gnu-} 前缀}\\
{GNU binutils} & {2.38}\\
{宿主 GCC / G++} & {11.4.0}\\
{Make / CMake / Ninja} & {4.3 / 3.22.1 / 1.10.1}\\
{Python / GTKWave} & {3.10.12 / 3.3.104}\\
\bottomrule\end{longtable}}


完整输出见\seqsplit{tool-versions.txt}。Ubuntu 软件包清单用于说明原环境的软件来源，新环境优先安装所需工具并核对关键版本。Mill 读取仓库 .mill-version 选择项目版本，启动器的安装来源与项目实际运行版本应分别识别；Scala、Chisel 和 firtool 则按本项目解析结果核对。\par

\subsection{跨环境复现条件}

环境部署涉及 Windows 上的工作树管理和 WSL 中的硬件构建。选择执行位置时，首先考虑文件所在的文件系统：Git 的索引检查需要访问大量文件，Scala、Chisel 和 C++ 构建也会产生大量中间产物。把工具放在相应文件系统的原生环境执行，可以减少跨环境文件访问，同时使路径、权限和工具行为更容易对应。\par


{\small\setlength{\tabcolsep}{3pt}\begin{longtable}{>{\raggedright\arraybackslash}p{0.27\linewidth}>{\raggedright\arraybackslash}p{0.3\linewidth}>{\raggedright\arraybackslash}p{0.35\linewidth}}
\caption{跨环境复现条件}\label{tab:1-8}\\
\toprule
\textbf{工作对象} & \textbf{执行位置} & \textbf{配套记录}\\\midrule\endfirsthead
\caption[]{跨环境复现条件（续）}\\
\toprule
\textbf{工作对象} & \textbf{执行位置} & \textbf{配套记录}\\\midrule\endhead
{Windows NTFS 工作树中的 Git 检查、提交与推送} & {Windows 原生 Git，使用 Windows 路径} & {分支、完整提交号、工作区差异及索引状态}\\
{Linux 或 WSL 原生目录中的 Git 操作} & {Linux 原生 Git} & {实际仓库根目录与提交身份}\\
{Chisel/Scala、硬件生成、RISC-V 编译、EMU/NEMU 和波形分析} & {Linux 或 WSL 原生文件系统} & {工具版本、构建参数、输入文件及输出目录}\\
{Windows 与 Linux 之间的源码、程序和证据交换} & {固定快照或独立检出，在接收端核验} & {完整提交号、未提交差异、依赖版本及关键文件 SHA-256}\\
\bottomrule\end{longtable}}


开始构建前，记录仓库根目录、提交号和实际输入；输出单独放置，避免与正在分析的旧记录混合。源码中存在未提交修改时，提交号应与差异快照共同描述输入。跨环境传输后，再核对用于运行的源码和二进制摘要，这样报告中的程序、运行日志和波形才能对应到同一次实验。\par


这一安排也受到后续工程实践的支持。2026-10-24 的 R4 调用入口发布检查中，一次由 WSL 访问 Windows 工作树的检查在约 300 秒后超时。核对确认未产生提交、没有残留进程或索引锁后，Windows 原生 Git 完成了剩余检查与提交。过程见原始发布检查和\href{https://github.com/cc77-builder/riscv/issues/2\#issuecomment-6376672788}{Issue 记录}。这是文件系统与工具选择的实际案例；两次执行条件不同，因此记录的是处理过程，没有据此计算 Git 操作的加速比例。\par

\subsection{10月25日验证的版本与产物}

T04运行绑定T17硬件生成来源\seqsplit{34bae6d4acceac939f958199dae53c14bc1f6183}，两次运行共同BIN的SHA-256为\seqsplit{d81437b44473cef2b0c5777330e13fd076567d4fc1a4296be35d971e51a4dcad}。生产EMU、FST模型、普通NEMU、flash及波形产物各自保留独立摘要。\par
本报告证据目录\seqsplit{evidence/stage1-current-20261025/}保存以下七份原件，其固定公开版本为2690c737：
\begin{itemize}
\item \href{https://github.com/cc77-builder/riscv/blob/2690c7370fb7c94111b4619bbac352e9be4ccc4d/team/evidence/stage1-current-20261025/production-run/emulator-output.log}{production-run/emulator-output.log}：英文hello、身份状态、用例结果和整程序统计。
\item \href{https://github.com/cc77-builder/riscv/blob/2690c7370fb7c94111b4619bbac352e9be4ccc4d/team/evidence/stage1-current-20261025/production-run/run.json}{production-run/run.json}：实际命令、输入前后摘要、日志摘要与退出状态。
\item \href{https://github.com/cc77-builder/riscv/blob/2690c7370fb7c94111b4619bbac352e9be4ccc4d/team/evidence/stage1-current-20261025/wave/wave.json}{wave/wave.json}：采样定义、信号清单、事件及选通信号产物摘要。
\item \href{https://github.com/cc77-builder/riscv/blob/2690c7370fb7c94111b4619bbac352e9be4ccc4d/team/evidence/stage1-current-20261025/acceptance.json}{acceptance.json}：软件输入等价、原输入核验、最终波形与工具检查汇总。
\item \href{https://github.com/cc77-builder/riscv/blob/2690c7370fb7c94111b4619bbac352e9be4ccc4d/team/evidence/stage1-current-20261025/RESULT.md}{RESULT.md}：两次运行、FST构建与完整产物的复现索引。
\item \href{https://github.com/cc77-builder/riscv/blob/2690c7370fb7c94111b4619bbac352e9be4ccc4d/team/evidence/stage1-current-20261025/trace-run/emulator-output.log}{trace-run/emulator-output.log}：独立FST模型的原始串口与终态。
\item \href{https://github.com/cc77-builder/riscv/blob/2690c7370fb7c94111b4619bbac352e9be4ccc4d/team/evidence/stage1-current-20261025/trace-run/run.json}{trace-run/run.json}：FST采集命令、运行输入与退出状态。
\end{itemize}
\subsection{技术资料索引}
{\small\noindent \href{https://github.com/cc77-builder/riscv/blob/347b1ae7eafb0f30c3b2cc1f8da8d9c8cdd3f98e/team/docs/environment.md}{环境部署详细方法}。\par}
\section{原始操作日志}
\subsection{10月25日生产模型原始输出}
以下列出该次构建或运行的记录。
\VerbatimInput[breaklines,breakanywhere,breaksymbolleft={},breaksymbolright={},fontsize=\footnotesize,baselinestretch=1]{evidence/rendered/hello-vadd-production-20261025.txt}
\subsection{10月25日FST模型原始输出}
以下列出该次构建或运行的记录。
\VerbatimInput[breaklines,breakanywhere,breaksymbolleft={},breaksymbolright={},fontsize=\footnotesize,baselinestretch=1]{evidence/rendered/hello-vadd-trace-20261025.txt}
\subsection{程序构建原始日志}
以下列出该次构建或运行的记录。
\VerbatimInput[breaklines,breakanywhere,breaksymbolleft={},breaksymbolright={},fontsize=\footnotesize,baselinestretch=1]{evidence/workload-build.log}
\subsection{历史hello与vadd运行摘要}
以下列出该次构建或运行的记录。
\VerbatimInput[breaklines,breakanywhere,breaksymbolleft={},breaksymbolright={},fontsize=\footnotesize,baselinestretch=1]{evidence/hello-vadd-summary.log}
\subsection{NEMU构建原始日志}
以下列出该次构建或运行的记录。
\VerbatimInput[breaklines,breakanywhere,breaksymbolleft={},breaksymbolright={},fontsize=\footnotesize,baselinestretch=1]{evidence/rendered/nemu-build.txt}
\subsection{EMU构建原始日志}
以下列出该次构建或运行的记录。
\VerbatimInput[breaklines,breakanywhere,breaksymbolleft={},breaksymbolright={},fontsize=\footnotesize,baselinestretch=1]{evidence/emu-build-20261009-204228.log}
\subsection{工具与编译链版本原始输出}
以下列出该次构建或运行的记录。
\VerbatimInput[breaklines,breakanywhere,breaksymbolleft={},breaksymbolright={},fontsize=\footnotesize,baselinestretch=1]{evidence/tool-versions.txt}
\subsection{固定子模块版本清单}
以下列出该次构建或运行的记录。
\VerbatimInput[breaklines,breakanywhere,breaksymbolleft={},breaksymbolright={},fontsize=\footnotesize,baselinestretch=1]{evidence/submodules-after-build.lock.txt}
\subsection{原始波形事件记录}
以下列出该次构建或运行的记录。
\VerbatimInput[breaklines,breakanywhere,breaksymbolleft={},breaksymbolright={},fontsize=\footnotesize,baselinestretch=1]{evidence/vadd-commit-events.csv}
\subsection{向量运算单元事件记录}
以下列出该次构建或运行的记录。
\VerbatimInput[breaklines,breakanywhere,breaksymbolleft={},breaksymbolright={},fontsize=\footnotesize,baselinestretch=1]{evidence/vector-fix-alu-events.csv}
